\begin{textsample}{POS Dim 3 – rn – Score 19.00 – rn/rn\_inf\_12.txt}  \label{ex:f3_pos_099}
A \textbf{brincadeira} é vista enquanto atividade, linguagem e modo das \textbf{crianças} se relacionarem com o real e, constitui, portanto, contexto de aprendizagem e desenvolvimento.
Logo, assume características que a distingue de outros modos de atividade ( VIGOTSKI, 2007 ; BROUGÉRE, 1998 ; LOPES, 2005 ), tais como : a liberdade, a regra e a imaginação ( fantasia ), a ação voluntária da \textbf{criança} – a decisão de iniciar e findar a participação – a possibilidade de câmbio de papéis e de retomadas infinitas, entre outras.
Desta maneira, nem sempre podem ser “ controladas pelos \textbf{adultos} ”, muito embora eles possam, enquanto responsáveis pelas \textbf{crianças}, exercer o papel de mediadores e participantes da \textbf{brincadeira}, no sentido de ampliar suas possibilidades como experiência da \textbf{criança} e, assim, no cotidiano das \textbf{instituições}, podem favorecer situações com que produzam repertórios, narrativas, significados e sentidos que podem, por sua vez, enriquecer as \textbf{brincadeiras} das \textbf{crianças}, pois : Todas as formas de \textbf{brincadeira} aprendidas pelas \textbf{crianças} são enriquecidas com o trabalho feito no conjunto das experiências por elas vividas nas outras dimensões, como : a da linguagem verbal e contação de histórias ; a dimensão das linguagens artísticas e também dos saberes que a \textbf{criança} vai construindo enquanto pensa o mundo social e o da natureza ; a dimensão do conhecimento de medidas, proporção, quantidades, e a importante dimensão da experiência de \textbf{cuidar} do outro e de si.
Conforme o olhar sensível dos professores acompanha e estimula as \textbf{crianças} a explorar eventos físicos ao \textbf{brincar}, a produzir desenhos, a desenvolver sua corporeidade e a realizar um conjunto valioso de aprendizagens, elas podem revolucionar o seu desenvolvimento, o que, por sua vez, alimenta os conteúdos das \textbf{brincadeiras} de faz de conta que criam ( \textbf{OLIVEIRA}, 2011, pp. 88-89 ).
Para Wallon ( 2005 ) a ludicidade é a \textbf{marca} primordial da \textbf{infância}, entendida, também, como período da vida em que todas as atividades humanas encontram -se em fase nascente.
Nessa perspectiva, atividades como andar, correr, pular podem assumir feições \textbf{lúdicas}, por sua repetição, pelo prazer que provocam.
Para o autor, os “ jogos ” são, inicialmente, “ puramente funcionais ”, depois assumem caráter de ficção, de aquisição e de fabricação ( Ibid, p. 73 ).
Enquanto os primeiros são experimentações e repetições de ações que produzem efeitos ( \textbf{sons}, movimentos ), os “ jogos de ficção ” consistem nas \textbf{brincadeiras} de faz-de-conta, que envolvem “ atividades, cuja interpretação é mais complexa ” ( Ibid, p. 74 ).
Os “ jogos de aquisição ” envolvem exploração do ambiente, dos objetos e seres, \textbf{curiosidade} e busca de compreensão.
Os “ jogos de fabricação ”, por sua vez, envolvem elaborações, intervenções sobre os objetos : “ reunir, combinar, modificar, transformar objetos e criar novos ” ( WALLON, 2005, p. 74 ).
O autor destaca que, nesses jogos, tanto as atividades de “ aquisição ”, como de “ ficção ” têm papel importante.
Desse modo, afirma que \textbf{brincar} é um modo fundamental de a \textbf{criança} se relacionar com o mundo, de explorá -lo, compreendê -lo, produzir sentidos próprios.
É possível pensar e organizar, dessa forma, modos de intervenção nas interações e \textbf{brincadeiras}, assumindo a perspectiva de que a \textbf{criança} aprende e se desenvolve na \textbf{brincadeira}, especialmente aquela voltada ao faz de conta.
Cabe ao professor possibilitar experiências que envolvam interações e \textbf{brincadeiras} que promovam o acesso e a apropriação de linguagens e conhecimentos ( DANTAS, 2016 ).
APROFUNDANDO O TEMA Interações e \textbf{Brincadeiras} Kishimoto ( 2010 ) orienta que as interações que compõem os eixos curriculares são assim explicitadas e caracterizadas : Interação com a \textbf{professora} - O \textbf{brincar} interativo com a \textbf{professora} é essencial para o conhecimento do mundo social e para dar maior riqueza, complexidade e qualidade às \textbf{brincadeiras}.
Especialmente para \textbf{bebês}, são essenciais ações \textbf{lúdicas} que envolvam turnos de falar ou gesticular, \textbf{esconder} e achar objetos.
Interação com as \textbf{crianças} - O \textbf{brincar} com outras \textbf{crianças} garante produção, conservação e recriação do repertório \textbf{lúdico} \textbf{infantil}.
Essa modalidade de cultura é conhecida como cultura \textbf{infantil} ou cultura \textbf{lúdica}.
Interação com os \textbf{brinquedos} e materiais - É essencial para o conhecimento do mundo dos objetos.
A diversidade de formas, \textbf{texturas}, cores, tamanhos, espessuras, cheiros e outras especificidades do objeto são importantes para a \textbf{criança} compreender esse mundo.
Interação entre \textbf{criança} e ambiente - A organização do ambiente pode facilitar ou dificultar a realização das \textbf{brincadeiras} e das interações entre as \textbf{crianças} e os \textbf{adultos}.
O ambiente físico reflete as concepções que a \textbf{instituição} assume para educar a \textbf{criança}.
Interações ( relações ) entre a \textbf{Instituição}, a família e a \textbf{criança} - A relação entre a \textbf{instituição} e a família possibilita o conhecimento e a inclusão, no projeto pedagógico, da cultura popular e dos \textbf{brinquedos} e \textbf{brincadeiras} que a \textbf{criança} conhece. ( Ibid ; pp. 02-03 ).
Destaca -se que a qualidade das interações entre professores e \textbf{crianças} pode facilitar ou dificultar o diálogo e as ações de reciprocidade.
Defende -se, portanto, que as professoras estabeleçam uma relação que envolva corpo e olhar numa perspectiva de igualdade.
Como, por exemplo, sentar na \textbf{roda} junto às \textbf{crianças} no lugar de observar de fora, e inclinar -se para ‘ olhar ’ nos olhos das \textbf{crianças}, no lugar de olhar de cima, impondo ação. “ Essa forma de interação possibilita a construção de uma cultura partilhada entre a \textbf{professora} e a \textbf{criança}, criando um fluxo positivo que potencializa a aprendizagem ” ( BRASIL, 2012, p. 16 ).

% matched lemmas: adulto, bebê, brincadeira, brincar, brinquedo, criança, cuidar, curiosidade, esconder, infantil, infância, instituição, lúdico, marca, oliveira, professora, roda, som, textura
\end{textsample}
